\Needspace{8\baselineskip}\originalchapter{线性稳定性与能量原理}\label{ch:energy}
\originalsection{从位移到力算符}
静态理想平衡上, 令位移 \ensuremath{\xi}, 质量、压力和磁场扰动为
\begin{equation}
 \rho_1=-\nabla\cdot(\rho\boldsymbol\xi),\quad
 p_1=-\boldsymbol\xi\cdot\nabla p-\gamma p\nabla\cdot\boldsymbol\xi,\quad
 \mathbf Q=\mathbf B_1=\nabla\times(\boldsymbol\xi\times\mathbf B).
\end{equation}
代线性动量得 $\rho\ddot{\boldsymbol\xi}=\mathcal F[\boldsymbol\xi]$, 其中
$\mathcal F=-\nabla p_1+\mu_0^{-1}(\nabla\times\mathbf Q)\times\mathbf B+\mathbf J\times\mathbf Q$.
给合适边界后积分分部使 $\int\boldsymbol\eta^*\cdot\mathcal F[\boldsymbol\xi]$ 对 \ensuremath{\xi}、\ensuremath{\eta} 共轭交换对称. 自伴性依赖静态平衡、理想条件与边界; 有平衡流或耗散不直接套用.

\originalsection{势能泛函与必要量词}
定义 $\delta W=-\tfrac12\operatorname{Re}\int\boldsymbol\xi^*\cdot\mathcal F[\boldsymbol\xi]\dd V$. 固定导体边界的流体项写为
\begin{equation}\label{eq:energy}
 \delta W_F=\frac12\operatorname{Re}\int\left[
 \frac{|\mathbf Q|^2}{\mu_0}-\boldsymbol\xi^*\cdot(\mathbf J\times\mathbf Q)
 +\gamma p|\nabla\cdot\boldsymbol\xi|^2
 +(\boldsymbol\xi\cdot\nabla p)(\nabla\cdot\boldsymbol\xi^*)\right]\dd V.
\end{equation}
自由界面还要加表面项与真空磁能 $\int_V|B_1|^2/(2\mu_0)$. 真空扰动满足 curl=div=0, 壁上法向扰动零, 等离子体界面满足随动磁通条件. 漏掉真空能会改变外 kink 阈值.

若\textbf{所有}可允许位移 \ensuremath{\delta}W\ensuremath{\ge}0则理想线性稳定; 找到一个 \ensuremath{\delta}W<0 的试探位移即足以证明不稳定. 找不到负值只表示所选试探空间未发现不稳定, 不能把几个系数的正定判成全MHD稳定. 对正规模 $\omega^2=\delta W/K$, $K=\tfrac12\int\rho|\xi|^2\dd V>0$; 负势能对应指数增长.

\originalsection{自由能怎样进入}
将式\eqref{eq:energy}改写可辨磁线弯曲、磁压缩、气体压缩的非负项, 以及压力梯度与曲率的耦合 $-2(\xi_\perp\cdot\nabla p)(\kappa\cdot\xi_\perp^*)$ 和平行电流耦合. 压力由高压向外驱动, 磁场曲率若向高压侧为坏曲率; 线弯曲和剪切可恢复. 这给 interchange、ballooning 与 current-driven kink 的物理区分. 平行位移可降低压缩能, 但在 $\mathbf k\cdot\mathbf B=0$ 的共振面其形式可能奇异, 必须检验可允许性而不能无条件设 div\ensuremath{\xi}=0.

\originalsection{有限谐波变分计算}
取满足轴心、壁及散度条件的基函数 $\xi=\sum_jc_j\xi_j$, 得 $\delta W=\tfrac12c^\dagger Wc$. 矩阵需 Hermitian, 实基可为实对称. 若要增长率还需质量矩阵 K 作广义特征值问题. 仅判断势能符号时最小 W 特征值足够, 基的缩放改变特征值大小但不改惯性指数. 增加基数可发现先前遗漏的负方向, 而数值求积收敛与基空间收敛是两种检查.
